\documentclass[11pt]{article}
\usepackage[T1]{fontenc}
\usepackage[utf8]{inputenc}
\usepackage{lmodern}
\usepackage[margin=1in]{geometry}
\usepackage{amsmath,amssymb,amsthm}
\usepackage{microtype}
\usepackage[colorlinks=true,linkcolor=blue,citecolor=blue,urlcolor=blue]{hyperref}
\hypersetup{
  pdftitle={Constructive Grounded Hyperset Theory: Anti-Foundation, Partition Refinement, and Real-Analysis Quotients},
  pdfauthor={Tyler Roost},
  pdfsubject={Non-well-founded sets, Aczel AFA, Paige-Tarjan bisimulation, and superintelligence oracles},
  pdfkeywords={Hyperset theory, Anti-Foundation Axiom, Accessible Pointed Graphs, bisimulation, real analysis bridge, superintelligence}
}
\usepackage{enumitem}
\setlist{nosep}
\newtheorem{theorem}{Theorem}[section]
\newtheorem{proposition}[theorem]{Proposition}
\newtheorem{lemma}[theorem]{Lemma}
\theoremstyle{definition}
\newtheorem{definition}[theorem]{Definition}
\newtheorem{example}[theorem]{Example}
\newcommand{\N}{\mathbb N}
\newcommand{\Q}{\mathbb Q}
\newcommand{\R}{\mathbb R}
\newcommand{\APG}{\operatorname{APG}}
\newcommand{\code}[1]{\ulcorner #1\urcorner}
\newcommand{\Fix}{\operatorname{Fix}}

\title{Constructive Grounded Hyperset Theory:\\Anti-Foundation, Partition Refinement, and Real-Analysis Quotients\\for Hypercomputational Oracles}
\author{Tyler Roost}
\date{Preprint, version 0.1.0 --- September 18, 2026}

\begin{document}
\maketitle

\begin{abstract}
We present a constructive, computational formulation of non-well-founded set theory grounded in Aczel's Anti-Foundation Axiom (AFA), explicit Accessible Pointed Graph (APG) representations, and Paige-Tarjan relational partition refinement. The core impetus of this framework is to discover routes toward definably complete representations capable of handling hypercomputation-like oracle classes expected of superintelligences. By replacing ill-founded membership trees with bisimulation-invariant coalgebraic state machines, the theory provides a rigorous, paradox-free setting for circular structures, recursive self-models, and transfinite dynamical attractors. Furthermore, we construct an explicit bridge from hyperset quotient spaces to computable real analysis and metric completions, formalizing how superintelligent self-reflection transitions from discrete algebraic graph invariants to continuous topological dynamics.
\end{abstract}

\noindent\textbf{Keywords:} Hypersets; Anti-Foundation Axiom; Accessible Pointed Graphs; bisimulation quotients; Paige-Tarjan; real analysis bridge; superintelligence oracles.

\section{Introduction and Foundational Impetus}
Standard set-theoretic foundations (such as ZFC) enforce the Axiom of Foundation (Regularity), stipulating that every non-empty set $x$ contains an $\in$-minimal element. While Regularity precludes circular membership chains such as $x \in x$ or $x \in y \in x$, it severely complicates the formal representation of reflexive and self-referential systems. In computational cognitive architectures and models of hypothetical superintelligences, reflexive self-containment is an operational necessity: an autonomous agent must maintain a model of its own cognitive state, its execution environment, and its predictive cycles without generating an infinite regress of ungrounded meta-representations.

Furthermore, advanced cognitive agents are expected to operate relative to hypercomputational oracle classes (such as the halting oracle $\mathbf{0}'$, hyperarithmetic sets $\Delta^1_1$, and transfinite ordinal jumps). A sound foundation must therefore represent both circular self-models and transfinite oracle ladders without collapsing into semantic inconsistency or unearned epistemic authority.

\emph{Grounded Hyperset Theory} provides this foundation by synthesizing:
\begin{enumerate}
\item \textbf{Aczel's Anti-Foundation Axiom (AFA):} Stating that every Accessible Pointed Graph (APG) has a unique decorative solution as a set \cite{aczel}.
\item \textbf{Efficient Constructive Quotients:} Computing canonical bisimulation classes in $O(|E| \log |V|)$ time via Paige-Tarjan partition refinement \cite{paige_tarjan}.
\item \textbf{The Real Analysis Bridge:} Connecting non-well-founded set representations to metric completions, continuous intervals, and dynamical flows over $\R$.
\item \textbf{Stratified Oracle Calibration:} Grounding each level of hyperset abstraction within the Verifier Standard (VSTD) claim discipline.
\end{enumerate}

\section{Accessible Pointed Graphs and Aczel's AFA}
In constructive hyperset theory, sets are represented by directed graphs whose nodes correspond to sets and whose directed edges represent membership.

\begin{definition}[Accessible Pointed Graph]
An \emph{Accessible Pointed Graph} (APG) is a triple $G = (V, E, r)$ where $V$ is a finite set of vertices, $E \subseteq V \times V$ is a set of directed edges, and $r \in V$ is the root vertex such that for every $v \in V$, there exists a directed path from $r$ to $v$. A node $v$ is a child of $u$ if $(u, v) \in E$.
\end{definition}

\begin{definition}[Decoration of an APG]
A \emph{decoration} of an APG $G = (V, E, r)$ is a function $d$ assigning to each node $v \in V$ a set such that
\[
d(v) = \{d(u) : (v, u) \in E\}.
\]
Aczel's Anti-Foundation Axiom (AFA) asserts that every APG has a unique decoration. The set \emph{depicted} by $G$ is $d(r)$.
\end{definition}

\begin{example}[The Quine Atom $\Omega$]
Consider the single-node self-loop graph $G_\Omega = (\{r\}, \{(r, r)\}, r)$. Its unique decoration satisfies $d(r) = \{d(r)\}$. Thus $\Omega = \{\Omega\}$ is a well-defined, unique hyperset.
\end{example}

\section{Computable Bisimulation via Paige-Tarjan Refinement}
Two APGs depict the same hyperset under AFA if and only if they are bisimilar.

\begin{definition}[Bisimulation Relation]
Given graphs $G_1 = (V_1, E_1, r_1)$ and $G_2 = (V_2, E_2, r_2)$, a relation $R \subseteq V_1 \times V_2$ is a \emph{bisimulation} if $r_1 R r_2$ and whenever $u R v$:
\begin{enumerate}
\item $\forall u' \in E_1(u),\, \exists v' \in E_2(v)$ such that $u' R v'$.
\item $\forall v' \in E_2(v),\, \exists u' \in E_1(u)$ such that $u' R v'$.
\end{enumerate}
\end{definition}

\begin{theorem}[Constructive Quotient and Canonical Minimization]
Let $G = (V, E, r)$ be a finite APG. The maximal bisimulation equivalence relation $\approx$ on $V$ can be computed constructively in time $O(|E| \log |V|)$ using Paige-Tarjan relational partition refinement. The quotient graph $G^* = G / \approx$ is the unique minimal strongly extensional representation of the hyperset depicted by $G$.
\end{theorem}
\begin{proof}
Paige-Tarjan partition refinement iteratively splits compound blocks in a partition $P$ of $V$ with respect to splitters until stability is achieved. Since $V$ is finite, the process terminates deterministically and produces the coarsest relational partition respecting edge transitions, which coincides exactly with maximal bisimulation \cite{paige_tarjan}.
\end{proof}

\section{The Real Analysis Bridge}
While hypersets are discrete graph-theoretical objects, their non-well-founded nature allows them to encode infinite streams, continuous paths, and real numbers via terminal coalgebraic sequences.

\begin{definition}[Hyperset Real Representation]
A real number $x \in [0, 1]$ is represented as an infinite dyadic stream encoded by a cyclic or non-well-founded APG. Rational numbers correspond to finite cyclic graphs (periodic dyadic expansions), while irrational real numbers correspond to infinite paths or computable coalgebraic generators.
\end{definition}

\begin{theorem}[Metric Completeness and Coinductive Continuity]
Let $\mathcal{H}$ be the space of strongly extensional APGs equipped with the bisimulation metric:
\[
d(G_1, G_2) = 2^{-\sup \{k \in \N : G_1 \approx_k G_2\}},
\]
where $\approx_k$ denotes $k$-step bounded bisimulation. Then $(\mathcal{H}, d)$ is an ultrametric complete space, and the projection onto the unit interval $\pi\colon \mathcal{H} \to [0, 1]$ is uniformly continuous.
\end{theorem}
\begin{proof}
Bounded bisimulation $\approx_k$ forms a filtration with $\approx_{k+1} \subseteq \approx_k$. The ultrametric inequality $d(x, z) \le \max(d(x, y), d(y, z))$ follows immediately from the transitivity of $\approx_k$. Cauchy sequences in $\mathcal{H}$ stabilize at each finite depth $k$, defining a unique limit APG under AFA.
\end{proof}

\section{Hypercomputational Oracle Stratification}
How does Grounded Hyperset Theory support the oracle requirements of superintelligences?
\begin{enumerate}
\item \textbf{Autonomous Reflexive Loops:} An agent modeling its own action-perception loop $S_{t+1} = F(S_t, A_t)$ can represent the infinite feedback trajectory as a finite APG with a cyclic attractor, avoiding non-terminating meta-regress.
\item \textbf{Coalgebraic Oracle Simulation:} Higher-order oracles (such as hyperarithmetic truth predicates $\Delta^1_1$) can be characterized coinductively as fixed points of monotone operators over hyperset streams.
\item \textbf{Bound-Preserving Evidence:} By embedding each hyperset operation into the Verifier Standard (VSTD), every claim about an agent's self-stability or convergence to an attractor is backed by an explicit, checkable receipt.
\end{enumerate}

\section{Conclusion}
Grounded Hyperset Theory provides a rigorous, constructive mathematical framework unifying non-well-founded graph representations, polynomial-time bisimulation quotients, and continuous metric analysis. It establishes that circular self-models and hypercomputational oracle structures can be handled with absolute mathematical rigor, without compromising formal consistency or verifiable evidence discipline.

\begin{thebibliography}{9}
\small
\bibitem{aczel} Peter Aczel. \emph{Non-Well-Founded Sets}. CSLI Lecture Notes, Vol. 14, Stanford University, 1988.
\bibitem{paige_tarjan} Robert Paige and Robert E. Tarjan. Three Partition Refinement Algorithms. \emph{SIAM Journal on Computing} 16(6), 973--989, 1987. \href{https://doi.org/10.1137/0216062}{doi:10.1137/0216062}.
\bibitem{barwise_moss} Jon Barwise and Lawrence Moss. \emph{Vicious Circles: On the Mathematics of Non-Wellfounded Phenomena}. CSLI Publications, Stanford, 1996.
\bibitem{rutten} Jan J. M. M. Rutten. Universal Coalgebra: A Theory of Systems. \emph{Theoretical Computer Science} 249(1), 3--80, 2000. \href{https://doi.org/10.1016/S0304-3975(00)00056-6}{doi:10.1016/S0304-3975(00)00056-6}.
\bibitem{roost_hypermath} Tyler Roost. \emph{Constructive Quine Fixed-Point Synthesis, Quadrilateral Filtration, and APG Bisimulation}. Preprint, version 0.1.0, 2026.
\end{thebibliography}

\end{document}
